% ============================================================
% Chapitre 2 — Analyse des besoins et spécification des exigences
% VisionGuard PFE — ENISo 2025-2026 — Yahya Ben Turkia
% ============================================================

\chapter{Analyse des besoins et spécification des exigences}

% --------------------------------------------------------
\section{Introduction}
% --------------------------------------------------------

Ce chapitre a pour objectif de formaliser le besoin industriel identifié au chapitre
précédent selon une approche d'ingénierie système légère. Cette démarche vise à délimiter
le système à concevoir, à identifier les acteurs qui interagissent avec lui et à structurer
les exigences fonctionnelles et techniques qui guideront les choix de conception.

Dans un premier temps, le système est situé par rapport à son environnement à travers un
diagramme de contexte et l'identification de ses principales interactions externes. Les cas
d'utilisation sont ensuite présentés afin de décrire les scénarios d'interaction entre les
acteurs et le système. Les exigences sont enfin formalisées dans des tableaux structurés,
destinés à servir de référence pour la conception, la réalisation et la validation du
prototype. Le chapitre se termine par une analyse préliminaire des risques liés à la
conception, à l'intégration et à la validation du système.

% --------------------------------------------------------
\section{Délimitation du système et analyse des interactions}
% --------------------------------------------------------

\subsection{Diagramme de contexte}

Le diagramme de contexte permet de délimiter les frontières du système d'inspection
automatisé VisionGuard par rapport à son environnement extérieur. Il met en évidence les
acteurs externes qui interagissent directement avec le système, ainsi que les principaux
flux physiques, informationnels et énergétiques échangés en conditions réelles
d'utilisation.

Dans ce projet, le système VisionGuard est considéré comme un dispositif d'inspection
semi-automatisé destiné au contrôle des pièces M00. Il reçoit une pièce positionnée par
l'opérateur, réalise l'acquisition des surfaces internes à inspecter, traite les images
acquises et fournit un résultat d'inspection exploitable par l'opérateur et le service
qualité.

Les acteurs externes identifiés sont les suivants :

\begin{itemize}
    \item \textbf{Opérateur} : positionne la pièce M00 sur le support, lance le cycle
    d'inspection et consulte les résultats.
    \item \textbf{Service qualité} : définit les critères d'acceptation et exploite les
    résultats d'inspection pour le suivi de conformité.
    \item \textbf{Pièce M00} : objet physique soumis à l'inspection.
    \item \textbf{Alimentation électrique} : fournit l'énergie nécessaire au fonctionnement
    des sous-systèmes.
    \item \textbf{Environnement atelier} : impose des contraintes d'encombrement, de
    poussière, de manipulation et de robustesse.
    \item \textbf{Maintenance} : impose l'accessibilité aux composants, la facilité de
    réglage et l'entretien du système.
\end{itemize}

\begin{figure}[H]
    \centering
    \includegraphics[width=0.85\textwidth]{contexte_visionguard.png}
    \caption{Diagramme de contexte du système VisionGuard.}
    \label{fig:contexte}
\end{figure}

% --------------------------------------------------------
\subsection{Analyse des interactions système-environnement}
\label{sec:interactions}

À partir du diagramme de contexte, les principales interactions entre VisionGuard et son
environnement peuvent être précisées. Cette analyse permet de transformer les échanges
identifiés en exigences fonctionnelles, techniques ou contraintes de conception. Elle
constitue ainsi une étape intermédiaire entre la délimitation du système et la formulation
du cahier des charges.

Le tableau suivant synthétise les interactions retenues.

\renewcommand{\arraystretch}{1.3}
\begin{longtable}{|p{1.5cm}|p{6cm}|p{4cm}|p{4cm}|}
\caption{Interactions entre le système VisionGuard et son environnement.}
\label{tab:interactions} \\
\hline
\rowcolor{gray!20}
\textbf{ID} & \textbf{Interaction} & \textbf{Acteur externe} & \textbf{Nature de l'échange} \\
\hline
\endfirsthead
\hline
\rowcolor{gray!20}
\textbf{ID} & \textbf{Interaction} & \textbf{Acteur externe} & \textbf{Nature de l'échange} \\
\hline
\endhead
\hline
\endfoot
\textbf{I-01} & Chargement et positionnement de la pièce M00 & Opérateur & Entrée physique \\ \hline
\textbf{I-02} & Lancement du cycle d'inspection & Opérateur & Commande utilisateur \\ \hline
\textbf{I-03} & Consultation du résultat d'inspection & Opérateur & Sortie informationnelle \\ \hline
\textbf{I-04} & Exploitation des résultats et images enregistrées & Service qualité & Sortie données / traçabilité \\ \hline
\textbf{I-05} & Présentation des surfaces internes à inspecter & Pièce M00 & Entrée physique \\ \hline
\textbf{I-06} & Alimentation du système & Alimentation électrique & Entrée énergétique \\ \hline
\textbf{I-07} & Contraintes d'encombrement, de poussière et de manipulation & Environnement atelier & Contrainte d'intégration \\ \hline
\textbf{I-08} & Accès aux composants, réglages et opérations d'entretien & Maintenance & Contrainte de maintenabilité \\ \hline
\end{longtable}

% --------------------------------------------------------
\section{Cas d'utilisation}
% --------------------------------------------------------

\subsection{Identification des acteurs}

Dans le cadre du système VisionGuard, deux acteurs principaux interagissent avec le système
lors de son utilisation opérationnelle.

L'opérateur constitue l'acteur principal. Il assure le chargement de la pièce, le lancement
du cycle d'inspection et la consultation du résultat final. Son intervention est limitée aux
actions de préparation, de commande et de supervision~; il n'intervient pas dans le
déroulement automatique de la séquence d'acquisition.

Le service qualité constitue un acteur secondaire. Il exploite les résultats d'inspection
enregistrés à des fins de suivi de conformité, d'analyse des défauts et de traçabilité. Son
interaction avec le système est principalement différée, après la réalisation du cycle
d'inspection.

La maintenance est considérée comme un acteur de contexte. Elle n'intervient pas dans le
cycle nominal d'inspection, mais impose des contraintes de conception relatives à
l'accessibilité, au réglage, au diagnostic et au remplacement des composants.

% --------------------------------------------------------
\subsection{Diagramme de cas d'utilisation}
\label{sec:usecase}

Le diagramme de cas d'utilisation représente les interactions fonctionnelles entre les
acteurs identifiés et le système VisionGuard. Il décrit les scénarios d'utilisation du
système sans préjuger des choix techniques de réalisation.

\begin{figure}[H]
    \centering
    \includegraphics[width=0.90\textwidth]{usecase_visionguard.png}
    \caption{Diagramme de cas d'utilisation du système VisionGuard.}
    \label{fig:usecase}
\end{figure}

% --------------------------------------------------------
\subsection{Description synthétique des cas d'utilisation}

Le tableau suivant présente une description synthétique de chaque cas d'utilisation
identifié, en précisant l'acteur concerné, le scénario nominal et le résultat attendu.

\renewcommand{\arraystretch}{1.3}
\small
\begin{longtable}{|p{1.5cm}|p{2.5cm}|p{5.5cm}|p{4cm}|}
\caption{Description synthétique des cas d'utilisation du système VisionGuard.}
\label{tab:usecase} \\
\hline
\rowcolor{gray!20}
\textbf{ID} & \textbf{Acteur} & \textbf{Scénario nominal} & \textbf{Résultat attendu} \\
\hline
\endfirsthead
\hline
\rowcolor{gray!20}
\textbf{ID} & \textbf{Acteur} & \textbf{Scénario nominal} & \textbf{Résultat attendu} \\
\hline
\endhead
\hline
\endfoot
\textbf{UC-01} & Opérateur
    & L'opérateur place la pièce M00 sur le support rotatif et vérifie son positionnement initial.
    & Pièce prête à être inspectée \\ \hline
\textbf{UC-02} & Opérateur
    & L'opérateur lance le cycle via l'interface utilisateur.
    & Le système démarre la séquence d'inspection. \\ \hline
\textbf{UC-03} & Opérateur
    & L'interface affiche l'état d'avancement de l'inspection en temps réel.
    & L'opérateur est informé de la progression du cycle. \\ \hline
\textbf{UC-04} & Opérateur
    & Le système affiche le résultat global PASS/FAIL et localise les zones défectueuses.
    & L'opérateur dispose du résultat d'inspection. \\ \hline
\textbf{UC-05} & Opérateur
    & L'opérateur interrompt le cycle en cours via l'interface.
    & Le système s'arrête proprement et conserve les données disponibles. \\ \hline
\textbf{UC-06} & Service qualité
    & Le service qualité consulte l'historique des inspections enregistrées.
    & Résultats disponibles avec horodatage et identifiant pièce. \\ \hline
\textbf{UC-07} & Service qualité
    & Le service qualité accède aux images associées aux défauts détectés.
    & Images disponibles pour analyse et décision. \\ \hline
\textbf{UC-08} & Service qualité
    & Le service qualité exporte les données d'inspection.
    & Un fichier de traçabilité est généré. \\ \hline
\end{longtable}

% --------------------------------------------------------
\section{Spécification des exigences}
% --------------------------------------------------------

\subsection{Exigences fonctionnelles}

Les exigences fonctionnelles décrivent ce que le système VisionGuard doit être capable de
réaliser. Elles sont directement dérivées des cas d'utilisation identifiés en
section~\ref{sec:usecase} et des interactions système-environnement définies en
section~\ref{sec:interactions}.

\renewcommand{\arraystretch}{1.3}
\begin{longtable}{|p{1.5cm}|p{10.5cm}|p{2.5cm}|}
\caption{Exigences fonctionnelles du système VisionGuard.}
\label{tab:ef} \\
\hline
\rowcolor{gray!20}
\textbf{ID} & \textbf{Exigence fonctionnelle} & \textbf{Source} \\
\hline
\endfirsthead
\hline
\rowcolor{gray!20}
\textbf{ID} & \textbf{Exigence fonctionnelle} & \textbf{Source} \\
\hline
\endhead
\hline
\endfoot
\textbf{EF-01} & Le système doit permettre le positionnement et l'indexation angulaire de la pièce M00 face aux veines à inspecter.
    & UC-01, I-01 \\ \hline
\textbf{EF-02} & Le système doit acquérir des images exploitables des parois internes des veines.
    & UC-02, I-05 \\ \hline
\textbf{EF-03} & Le système doit permettre l'observation des deux parois internes de chaque veine, soit par orientation contrôlée de la caméra, soit par réglage angulaire validé lors de la phase prototype.
    & UC-02, I-05 \\ \hline
\textbf{EF-04} & Le système doit permettre la détection automatique des défauts de type pore ou piqûre d'un diamètre minimal visé de 0,8~mm.
    & UC-02, I-05 \\ \hline
\textbf{EF-05} & Le système doit afficher un résultat global PASS/FAIL à l'issue de chaque cycle d'inspection.
    & UC-04 \\ \hline
\textbf{EF-06} & Le système doit localiser les veines ou zones présentant des défauts détectés.
    & UC-04, UC-07 \\ \hline
\textbf{EF-07} & Le système doit enregistrer les images acquises avec horodatage et identifiant pièce.
    & UC-06, I-04 \\ \hline
\textbf{EF-08} & Le système doit permettre l'export des données de traçabilité.
    & UC-08 \\ \hline
\textbf{EF-09} & Le système doit permettre l'arrêt propre du cycle en cours de fonctionnement.
    & UC-05 \\ \hline
\end{longtable}

% --------------------------------------------------------
\subsection{Exigences non fonctionnelles}

Les exigences non fonctionnelles définissent les contraintes de performance, d'intégration,
de robustesse et d'environnement auxquelles le système VisionGuard doit se conformer,
indépendamment des fonctions qu'il assure.

\renewcommand{\arraystretch}{1.3}
\small
\begin{longtable}{|p{1.5cm}|p{3cm}|p{3cm}|p{3.5cm}|p{2.5cm}|}
\caption{Exigences non fonctionnelles du système VisionGuard.}
\label{tab:en} \\
\hline
\rowcolor{gray!20}
\textbf{ID} & \textbf{Exigence} & \textbf{Critère de mesure} & \textbf{Niveau exigé} & \textbf{Source} \\
\hline
\endfirsthead
\hline
\rowcolor{gray!20}
\textbf{ID} & \textbf{Exigence} & \textbf{Critère de mesure} & \textbf{Niveau exigé} & \textbf{Source} \\
\hline
\endhead
\hline
\endfoot
\textbf{EN-01} & Temps de cycle d'inspection complet
    & Durée mesurée par pièce
    & $\leq$ 3 minutes
    & Objectif projet \\ \hline
\textbf{EN-02} & Répétabilité de l'indexation angulaire
    & Écart angulaire mesuré sur indexations répétées
    & Écart inférieur à $\pm$1° sur l'ensemble des positions
    & Besoin d'acquisition \\ \hline
\textbf{EN-03} & Résolution spatiale de détection
    & Taille minimale du défaut détectable
    & Défaut visé $\geq$ 0,8~mm de diamètre
    & Besoin qualité \\ \hline
\textbf{EN-04} & Encombrement du système
    & Surface au sol occupée
    & $\leq$ 1~m²
    & Contrainte atelier \\ \hline
\textbf{EN-05} & Alimentation électrique
    & Type de source utilisée
    & 220~V AC secteur standard avec conversions adaptées
    & Contrainte d'intégration \\ \hline
\textbf{EN-06} & Résistance à l'environnement atelier
    & Fonctionnement en présence de poussière céramique et de manipulations répétées
    & Composants protégés et montage robuste
    & Environnement atelier \\ \hline
\textbf{EN-07} & Maintenabilité
    & Accessibilité aux composants et aux réglages
    & Accès possible sans outillage spécifique lourd
    & Maintenance \\ \hline
\textbf{EN-08} & Stabilité mécanique pendant acquisition
    & Netteté des images après indexation
    & Image exploitable sans flou significatif pour la détection du défaut minimal visé
    & Besoin vision \\ \hline
\textbf{EN-09} & Fiabilité de communication entre modules
    & Continuité des échanges pendant un cycle complet
    & Aucune interruption ni commande non exécutée observée lors des tests
    & Besoin d'intégration \\ \hline
\textbf{EN-10} & Sécurité opérateur
    & Protection contre les éléments en mouvement et les risques électriques
    & Respect des pratiques de sécurité en atelier
    & Contrainte sécurité \\ \hline
\end{longtable}

% --------------------------------------------------------
\subsection{Critères de validation associés}

Les critères de validation permettent de vérifier, lors de la phase de tests, que les
exigences formulées sont effectivement satisfaites par le système réalisé. Ils constituent
le lien direct entre la spécification des exigences et les résultats de validation qui
seront présentés au chapitre~5.

\renewcommand{\arraystretch}{1.3}
\small
\begin{longtable}{|p{1.5cm}|p{7cm}|p{7cm}|}
\caption{Critères de validation associés aux exigences.}
\label{tab:validation} \\
\hline
\rowcolor{gray!20}
\textbf{ID} & \textbf{Critère de validation} & \textbf{Méthode de vérification} \\
\hline
\endfirsthead
\hline
\rowcolor{gray!20}
\textbf{ID} & \textbf{Critère de validation} & \textbf{Méthode de vérification} \\
\hline
\endhead
\hline
\endfoot
\textbf{EF-01} & Indexation correcte sur les positions correspondant aux veines
    & Mesure ou observation répétée de l'alignement sur plusieurs cycles \\ \hline
\textbf{EF-02} & Images nettes et exploitables des surfaces inspectées
    & Inspection visuelle des images acquises et contrôle de la netteté \\ \hline
\textbf{EF-03} & Couverture des deux parois internes de chaque veine
    & Vérification de la couverture optique sur une pièce complète ou sur un échantillon représentatif \\ \hline
\textbf{EF-04} & Détection de défauts de référence de taille compatible avec le seuil visé
    & Tests sur pièces non conformes connues ou défauts de référence disponibles \\ \hline
\textbf{EF-05} & Affichage du résultat PASS/FAIL en fin de cycle
    & Test fonctionnel sur pièce conforme et pièce non conforme \\ \hline
\textbf{EF-06} & Localisation correcte des veines ou zones défectueuses
    & Comparaison entre résultat système et inspection manuelle experte \\ \hline
\textbf{EF-07} & Enregistrement des images avec horodatage et identifiant pièce
    & Vérification des fichiers générés après cycle \\ \hline
\textbf{EF-08} & Export fonctionnel des données de traçabilité
    & Test d'export et vérification du contenu du fichier généré \\ \hline
\textbf{EF-09} & Arrêt propre sans perte des données déjà acquises
    & Test d'interruption en cours de cycle \\ \hline
\textbf{EN-01} & Temps de cycle mesuré inférieur ou égal à 3 minutes
    & Chronométrage sur plusieurs cycles d'inspection \\ \hline
\textbf{EN-02} & Écart angulaire inférieur à $\pm$1° sur indexations répétées
    & Mesure de l'écart de position sur plusieurs indexations successives \\ \hline
\textbf{EN-03} & Résolution spatiale compatible avec le défaut minimal visé
    & Test sur défaut réel ou étalon de dimension connue \\ \hline
\textbf{EN-04} & Encombrement inférieur ou égal à 1~m²
    & Mesure physique du système assemblé \\ \hline
\textbf{EN-08} & Absence de flou significatif sur les images acquises
    & Évaluation de la netteté sur une série d'images consécutives \\ \hline
\textbf{EN-09} & Communication stable sur une séquence complète
    & Test de plusieurs cycles sans erreur de communication observée \\ \hline
\end{longtable}

% --------------------------------------------------------
\section{Décomposition fonctionnelle du système}
% --------------------------------------------------------

\subsection{Sous-systèmes principaux}
\label{sec:decomposition}

La décomposition fonctionnelle du système VisionGuard permet d'organiser le système en
ensembles cohérents de fonctions. Elle ne décrit pas encore en détail les choix techniques
retenus, qui seront présentés au chapitre suivant, mais elle permet d'identifier les rôles
principaux nécessaires à la réalisation de l'inspection automatisée.

Le système est structuré autour de trois sous-systèmes principaux.

Le \textbf{sous-système mécanique} assure le positionnement relatif entre la pièce M00 et
le dispositif d'acquisition. Il permet l'indexation angulaire de la pièce afin de présenter
successivement les veines devant l'axe d'inspection. Il doit également garantir la
stabilité de l'ensemble pendant l'acquisition des images.

Le \textbf{sous-système vision} assure l'observation et l'acquisition des surfaces internes
à contrôler. Il regroupe les fonctions liées à la caméra, à l'objectif, à l'éclairage et à
la qualité de l'image. Son rôle principal est de fournir des images suffisamment nettes,
contrastées et reproductibles pour permettre l'analyse automatique des défauts.

Le \textbf{sous-système contrôle, traitement et décision} assure la coordination globale du
système. Il pilote la séquence d'inspection, synchronise les mouvements et les acquisitions,
gère le stockage des images, transmet les données au module d'analyse et présente le
résultat final à l'utilisateur. Il inclut également la fonction de traçabilité, nécessaire
au suivi des résultats d'inspection.

\begin{figure}[H]
    \centering
    \includegraphics[width=0.95\textwidth]{decomposition_visionguard.png}
    \caption{Schéma de décomposition fonctionnelle du système VisionGuard.}
    \label{fig:decomposition}
\end{figure}

% --------------------------------------------------------
\subsection{Cycle opérationnel d'inspection}

Le cycle opérationnel d'inspection décrit la séquence de fonctionnement du système
VisionGuard depuis le chargement de la pièce jusqu'à l'affichage du résultat final. Ce
cycle se déroule de manière semi-automatisée~: l'opérateur intervient pour préparer et
lancer l'inspection, tandis que la séquence d'acquisition, de traitement et d'enregistrement
est exécutée par le système.

Le cycle nominal se déroule comme suit~:

\begin{enumerate}
    \item L'opérateur place la pièce M00 sur le support rotatif et lance le cycle via
    l'interface utilisateur.
    \item Le système initialise la position de référence et aligne la pièce face à la
    première veine à inspecter.
    \item Le système acquiert une image de la première paroi interne selon l'orientation
    optique définie.
    \item L'orientation optique est ajustée afin d'acquérir l'image de la seconde paroi
    interne de la veine.
    \item Le système indexe la pièce vers la veine suivante selon un pas angulaire
    correspondant à la répartition des veines.
    \item Les étapes d'acquisition et d'indexation sont répétées pour couvrir l'ensemble
    des veines à inspecter.
    \item Les images acquises sont transmises au module de classification par intelligence
    artificielle pour analyse.
    \item Le résultat global PASS/FAIL est affiché à l'opérateur, accompagné de la
    localisation des zones présentant des défauts détectés le cas échéant.
    \item Les images et le résultat sont enregistrés avec horodatage et identifiant pièce
    pour assurer la traçabilité.
\end{enumerate}

\begin{figure}[H]
    \centering
    \includegraphics[width=\textwidth]{sequence_visionguard.png}
    \caption{Diagramme de séquence du cycle opérationnel d'inspection.}
    \label{fig:sequence}
\end{figure}

% --------------------------------------------------------
\section{Analyse préliminaire des risques}
% --------------------------------------------------------

L'analyse préliminaire des risques a pour objectif d'identifier les principales défaillances
potentielles du système et de proposer des actions préventives ou correctives avant la phase
de réalisation. Cette analyse s'appuie sur une évaluation de la criticité, de la probabilité
d'occurrence et du niveau de risque résultant.

Le tableau suivant présente une synthèse des risques identifiés lors de la phase de
conception, classés selon leur niveau de criticité.

\renewcommand{\arraystretch}{1.3}
\footnotesize
\begin{longtable}{|p{2.8cm}|p{2.8cm}|c|c|p{1.8cm}|p{7cm}|}
\caption{Analyse préliminaire des risques du système VisionGuard.}
\label{tab:risques} \\
\hline
\rowcolor{gray!20}
\textbf{Risque identifié} & \textbf{Cause probable} &
\textbf{C} & \textbf{P} &
\textbf{Niveau} &
\textbf{Action préventive / corrective} \\
\hline
\endfirsthead
\hline
\rowcolor{gray!20}
\textbf{Risque identifié} & \textbf{Cause probable} &
\textbf{C} & \textbf{P} &
\textbf{Niveau} &
\textbf{Action préventive / corrective} \\
\hline
\endhead
\hline
\endfoot
Images floues ou mal cadrées
    & Vibrations, mauvais réglage optique, profondeur de champ insuffisante
    & 5 & 3 & 15 — Élevé
    & Calibration optique, stabilisation mécanique, pause avant capture \\ \hline
Défaut de positionnement angulaire
    & Jeu mécanique, imprécision d'indexation
    & 4 & 3 & 12 — Moyen
    & Réduction mécanique, vérification des tolérances, tests de répétabilité \\ \hline
Éclairage insuffisant ou non uniforme
    & Positionnement inadéquat des LEDs, puissance trop faible
    & 5 & 4 & 20 — Élevé
    & Tests en conditions réelles, ajustement orientation et intensité \\ \hline
Dataset insuffisant ou déséquilibré
    & Nombre limité de pièces non conformes
    & 5 & 4 & 20 — Élevé
    & Collecte progressive, augmentation des données, validation indépendante \\ \hline
Modèle IA sous-performant
    & Dataset non représentatif, sur-apprentissage
    & 5 & 4 & 20 — Élevé
    & Validation croisée, séparation apprentissage/test, évaluation indépendante \\ \hline
Temps de cycle dépassé
    & Acquisition lente, traitement IA lent
    & 3 & 3 & 9 — Moyen
    & Optimisation séquence, réduction temps morts, parallélisation \\ \hline
Panne composant électronique
    & Défaillance caméra, moteur, driver
    & 4 & 2 & 8 — Faible
    & Documentation câblage, rechanges identifiés, tests unitaires \\ \hline
Accès optique insuffisant
    & Géométrie chevron, angle d'observation limité
    & 5 & 3 & 15 — Élevé
    & Tests optiques préalables, ajustement angle caméra \\ \hline
Vibrations pendant acquisition
    & Rigidité insuffisante, mouvement non stabilisé
    & 3 & 3 & 9 — Moyen
    & Renforcement châssis, pause stabilisation, contrôle fixation caméra \\ \hline
Perte de communication
    & Déconnexion USB/Ethernet, erreur protocole
    & 4 & 3 & 12 — Moyen
    & Protocole simple, accusés de réception, tests séquence complète \\ \hline
\multicolumn{6}{|l|}{\textit{C = Criticité (1--5) \quad P = Probabilité (1--5) \quad Niveau = C $\times$ P \quad Faible~: $\leq$8 \quad Moyen~: 9--12 \quad Élevé~: $\geq$13}} \\
\hline
\end{longtable}
% --------------------------------------------------------
\section{Conclusion}
% --------------------------------------------------------

Ce chapitre a formalisé le besoin industriel selon une approche d'ingénierie système légère.
Le diagramme de contexte a permis de délimiter les frontières du système VisionGuard et
d'identifier les acteurs externes en interaction avec lui. L'analyse des interactions a
structuré les échanges entre le système et son environnement, servant de base à la
spécification des exigences.

Les cas d'utilisation ont décrit les principaux scénarios d'interaction opérationnelle,
tandis que les exigences fonctionnelles et non fonctionnelles ont traduit le besoin en
critères mesurables et vérifiables. Les critères de validation associés établissent le lien
direct entre les exigences définies dans ce chapitre et les essais de validation qui seront
présentés au chapitre~5.

La décomposition fonctionnelle a organisé le système en trois sous-systèmes principaux~:
mécanique, vision, et contrôle-traitement-décision. Le cycle opérationnel d'inspection a
ensuite décrit la séquence générale de fonctionnement, depuis le chargement de la pièce
jusqu'à l'affichage et l'enregistrement du résultat. Enfin, l'analyse préliminaire des
risques a identifié les points critiques nécessitant une attention particulière, notamment
la qualité d'acquisition d'image, la couverture optique des parois internes, la constitution
du dataset et la performance du module IA.

Le chapitre suivant présentera les choix de conception détaillés pour chacun des
sous-systèmes, en justifiant les solutions techniques retenues au regard des exigences
formulées dans ce chapitre.